\documentclass[11pt]{article}

\usepackage[a4paper,margin=1in]{geometry}
\usepackage{amsmath,amssymb,amsthm,mathtools,bm}
\usepackage{booktabs,tabularx,array}
\usepackage{enumitem}
\usepackage{graphicx}
\usepackage{float}
\usepackage{microtype}
\usepackage[hidelinks]{hyperref}

\newtheorem{assumption}{Assumption}[section]
\newtheorem{definition}[assumption]{Definition}
\newtheorem{proposition}[assumption]{Proposition}
\newtheorem{theorem}[assumption]{Theorem}
\newtheorem{lemma}[assumption]{Lemma}
\newtheorem{corollary}[assumption]{Corollary}
\theoremstyle{remark}
\newtheorem{remark}[assumption]{Remark}
\newtheorem{example}[assumption]{Example}

\newcommand{\E}{\mathbb E}
\newcommand{\Pp}{\mathbb P}
\newcommand{\R}{\mathbb R}
\newcommand{\cL}{\mathcal L}
\newcommand{\cH}{\mathcal H_T}
\newcommand{\cX}{\mathbb X_T}
\newcommand{\HS}{\mathcal S_2}
\newcommand{\cW}{\mathcal W}
\newcommand{\cZ}{\mathcal Z}
\newcommand{\N}{\mathcal N}
\newcommand{\Ph}{\Phi}
\newcommand{\ph}{\phi}
\newcommand{\ip}[2]{\left\langle #1,#2\right\rangle}
\newcommand{\norm}[1]{\left\lVert #1\right\rVert}
\newcommand{\abs}[1]{\left\lvert #1\right\rvert}
\newcommand{\tr}{\operatorname{tr}}
\newcommand{\Law}{\mathcal L}
\newcommand{\op}{\mathrm{o}_{\Pp}}
\newcommand{\CG}{C_G}

\title{Work Package 6: Operator Sensitivity and Causal Response Geometry\\
\large Semigroup derivatives, Gaussian law response, Wasserstein stability, and optimal local intervention directions}
\author{}
\date{7 July 2026}

\begin{document}
\maketitle

\begin{abstract}
This note develops the operator-sensitivity layer of a causal linear-SPDE programme. Two perturbation regimes are separated. First, for a generator $A$ and a bounded, possibly non-commuting direction $B\in\cL(H)$, the perturbed semigroup generated by $A+\varepsilon B$ is Fr\'echet differentiable in operator norm, with Duhamel derivative and a quadratic remainder. This expansion is transferred to the stochastic mild solution in $L^2(\Omega;L^2(0,T;H))$, then to Gaussian field-functional means, variances, CDFs, tails, and quantiles. A common-noise coupling gives a uniform local $W_2$ bound and a tangent-law result. Second, the paper's primary intervention $aA_0$ is treated separately: its formal direction $A_0$ is unbounded and therefore is not covered by the bounded-perturbation theorem. Analytic-semigroup calculus yields the uniform derivative $\partial_aS_a(t)=tA_0S_a(t)$ and recovers the scalar formulas used by the preceding inference packages. Finally, the operator derivative is interpreted as a response covector on $\cL(H)$ and, after restriction to Hilbert-Schmidt directions, as a genuine Riesz gradient. Explicit mean, variance, and CDF gradients yield the optimal local intervention direction. A non-commuting finite-dimensional experiment verifies the analytic derivative, the quadratic law remainder, and the Wasserstein envelope.
\end{abstract}

\section{Scope and the two perturbation regimes}

Let $H$ and $U$ be real separable Hilbert spaces. The positive core uses a linear stochastic evolution equation
\begin{equation}\label{eq:base-spde}
 dU_t=AU_t\,dt+\sigma G\,dW_t,
 \qquad U_0=u_0,
\end{equation}
where $W$ is a cylindrical Wiener process on $U$ and $G\in\cL(U,H)$.

Work Package 6 studies the local effect of changing the post-intervention generator. Two regimes must be distinguished.

\begin{enumerate}[label=(\roman*)]
\item A bounded direction,
\begin{equation}\label{eq:bounded-pert}
 A_{\varepsilon,B}=A+\varepsilon B,
 \qquad B\in\cL(H),
\end{equation}
which may be non-diagonal and non-commuting.
\item A scalar rescaling,
\begin{equation}\label{eq:scaling-pert}
 A_a=aA_0,
\end{equation}
which is the main diffusivity intervention in the paper. Its formal direction is $A_0$, generally unbounded.
\end{enumerate}

\begin{remark}[Why the regimes cannot be merged]
The bounded perturbation theorem does not apply to $B=A_0$ unless $A_0$ is bounded, which it is not for parabolic differential operators. The scalar-rescaling family nevertheless has a strong derivative because analytic smoothing controls $tA_0S_a(t)$. Treating these as separate theorems avoids an invalid use of bounded perturbation theory.
\end{remark}

\section{Standing assumptions and target topology}

\begin{assumption}[Exponentially stable semigroup]\label{ass:semigroup}
The operator $A:D(A)\subset H\to H$ generates a strongly continuous semigroup $S=(S(t))_{t\ge0}$ satisfying
\begin{equation}\label{eq:stable}
 \norm{S(t)}_{\cL(H)}\le Me^{-\omega t},
 \qquad t\ge0,
\end{equation}
for some $M\ge1$ and $\omega>0$.
\end{assumption}

\begin{assumption}[Admissible additive noise]\label{ass:noise}
The operator $G\in\cL(U,H)$ satisfies
\begin{equation}\label{eq:admissible}
 \int_0^T\norm{S(t)G}_{\HS(U,H)}^2\,dt<\infty.
\end{equation}
\end{assumption}

The mild solution of \eqref{eq:base-spde} is
\begin{equation}\label{eq:mild}
 U_t=S(t)u_0+\sigma\int_0^tS(t-r)G\,dW_r.
\end{equation}
Set
\begin{equation}\label{eq:spaces}
 \cH:=L^2(0,T;H),
 \qquad
 \cX:=L^2(\Omega;\cH).
\end{equation}
For $w\in\cH$, define
\begin{equation}\label{eq:functional}
 \Psi_w(X):=\int_0^T\ip{w_t}{X_t}_H\,dt.
\end{equation}
Then
\begin{equation}\label{eq:functional-cont}
 \E\abs{\Psi_w(X)}^2
 \le
 \norm{w}_{\cH}^2\norm{X}_{\cX}^2.
\end{equation}

\begin{remark}[Why the integrated topology is sufficient]
The field-functional law depends continuously on the path in $\cX$. A derivative in $L^2(\Omega;C([0,T];H))$ would require stronger maximal-regularity assumptions and is not needed for the target. The topology in \eqref{eq:spaces} is therefore both sufficient and analytically economical.
\end{remark}

\section{Bounded operator perturbations}

Fix $R<\infty$ and let
\begin{equation}\label{eq:Bball}
 \mathbb B_R:=\{B\in\cL(H):\norm{B}\le R\}.
\end{equation}
Choose $\varepsilon_0>0$ such that
\begin{equation}\label{eq:epsilon}
 M\varepsilon_0R<\omega,
 \qquad
 \omega_*:=\omega-M\varepsilon_0R>0.
\end{equation}
Let $S_{\varepsilon,B}$ be the semigroup generated by $A+\varepsilon B$.

\begin{lemma}[Uniform bounded-perturbation estimate]\label{lem:pert-bound}
Under Assumption~\ref{ass:semigroup},
\begin{equation}\label{eq:pert-stability}
 \norm{S_{\varepsilon,B}(t)}
 \le
 Me^{-\omega_*t}
\end{equation}
for $B\in\mathbb B_R$, $\abs{\varepsilon}\le\varepsilon_0$, and $t\ge0$.
\end{lemma}

\begin{proof}
The bounded perturbation theorem and the Dyson--Phillips estimate give
$\norm{S_{\varepsilon,B}(t)}\le M\exp\{(-\omega+M\abs{\varepsilon}\norm{B})t\}$. Use \eqref{eq:epsilon}.
\end{proof}

Define the Duhamel derivative
\begin{equation}\label{eq:DS}
 \mathcal D_BS(t)
 :=
 \int_0^tS(t-s)BS(s)\,ds.
\end{equation}

\begin{theorem}[Semigroup Fr\'echet derivative in operator norm]\label{thm:semigroup-derivative}
Under Assumption~\ref{ass:semigroup}, uniformly over $B\in\mathbb B_R$,
\begin{equation}\label{eq:semigroup-expansion}
 S_{\varepsilon,B}(t)
 =
 S(t)+\varepsilon\mathcal D_BS(t)
 +\mathcal R_{\varepsilon,B}(t),
\end{equation}
where
\begin{equation}\label{eq:semigroup-remainder}
 \sup_{0\le t\le T}
 \norm{\mathcal R_{\varepsilon,B}(t)}
 \le
 C_T\varepsilon^2\norm{B}^2.
\end{equation}
Consequently,
\begin{equation}\label{eq:frechet-S}
 B\longmapsto\mathcal D_BS(\cdot)
\end{equation}
is a bounded linear map from $\cL(H)$ into $C([0,T];\cL(H))$.
\end{theorem}

\begin{proof}
The variation-of-constants identity in the orientation needed for the noise estimate is
\begin{equation}\label{eq:duhamel}
 S_{\varepsilon,B}(t)-S(t)
 =
 \varepsilon\int_0^t
 S_{\varepsilon,B}(t-s)BS(s)\,ds.
\end{equation}
Subtract $\varepsilon\mathcal D_BS(t)$ and apply \eqref{eq:duhamel} once more to the first semigroup factor. This gives
\begin{align}
 \mathcal R_{\varepsilon,B}(t)
 &=
 \varepsilon^2
 \int_0^t\int_0^{t-s}
 S_{\varepsilon,B}(t-s-r)BS(r)BS(s)
 \,dr\,ds.
 \label{eq:R-double}
\end{align}
The bounds \eqref{eq:stable} and \eqref{eq:pert-stability} imply
$\norm{\mathcal R_{\varepsilon,B}(t)}\le C\varepsilon^2\norm{B}^2t^2$ on $[0,T]$. Linearity is immediate from \eqref{eq:DS}, and boundedness follows from
\[
 \norm{\mathcal D_BS(t)}
 \le M^2t e^{-\omega t}\norm{B}.
\]
\end{proof}

\begin{proposition}[Noise-transfer derivative]\label{prop:noise-transfer}
Under Assumptions~\ref{ass:semigroup}--\ref{ass:noise},
\begin{align}
 &\int_0^T\norm{\mathcal D_BS(t)G}_{\HS}^2dt
 \le C_T\norm{B}^2,
 \label{eq:DSG}\\
 &\int_0^T
 \norm{\mathcal R_{\varepsilon,B}(t)G}_{\HS}^2dt
 \le C_T\varepsilon^4\norm{B}^4.
 \label{eq:RSG}
\end{align}
\end{proposition}

\begin{proof}
Let $q(t)=\norm{S(t)G}_{\HS}$, which belongs to $L^2(0,T)$ by \eqref{eq:admissible}. Formula \eqref{eq:DS} yields
\[
 \norm{\mathcal D_BS(t)G}_{\HS}
 \le M\norm{B}\int_0^te^{-\omega(t-s)}q(s)\,ds.
\]
Young's convolution inequality proves \eqref{eq:DSG}. In \eqref{eq:R-double}, keep $S(s)G$ as the rightmost factor. The remaining inner integration contributes a bounded kernel proportional to $t-s$ on $[0,T]$. A second application of Young's inequality proves \eqref{eq:RSG}.
\end{proof}

\begin{remark}[No Hilbert-Schmidt assumption on $G$]
The proposition uses smoothing by $S(s)$ before $G$ enters the Hilbert-Schmidt norm. It therefore permits cylindrical forcing, including the one-dimensional heat equation with $G=I$.
\end{remark}

\section{Derivative of the stochastic mild solution}

For $B\in\mathbb B_R$, let
\begin{equation}\label{eq:pert-mild}
 U_t^{\varepsilon,B}
 :=
 S_{\varepsilon,B}(t)u_0
 +\sigma\int_0^tS_{\varepsilon,B}(t-r)G\,dW_r.
\end{equation}
Define
\begin{equation}\label{eq:dotU}
 \dot U_t[B]
 :=
 \mathcal D_BS(t)u_0
 +\sigma\int_0^t\mathcal D_BS(t-r)G\,dW_r.
\end{equation}

\begin{theorem}[Fr\'echet derivative of the counterfactual field]\label{thm:mild-derivative}
Under Assumptions~\ref{ass:semigroup}--\ref{ass:noise}, the perturbation map is Fr\'echet differentiable at zero from $\cL(H)$ into $\cX$. More precisely,
\begin{equation}\label{eq:mild-remainder}
 \norm{U^{\varepsilon,B}-U-\varepsilon\dot U[B]}_{\cX}
 \le C_T\varepsilon^2\norm{B}^2.
\end{equation}
The derivative has the equivalent mild response form
\begin{equation}\label{eq:response-mild}
 \dot U_t[B]
 =
 \int_0^tS(t-s)BU_s\,ds
\end{equation}
in $\cX$. Thus it solves
\begin{equation}\label{eq:response-equation}
 d\dot U_t=A\dot U_t\,dt+BU_t\,dt,
 \qquad
 \dot U_0=0,
\end{equation}
in the mild sense.
\end{theorem}

\begin{proof}
For the deterministic component, Theorem~\ref{thm:semigroup-derivative} gives
\[
 \int_0^T\norm{\mathcal R_{\varepsilon,B}(t)u_0}^2dt
 \le C_T\varepsilon^4\norm{B}^4\norm{u_0}^2.
\]
For the stochastic component, It\^o isometry and Fubini give
\begin{align*}
 &\E\int_0^T
 \norm{
 \sigma\int_0^t\mathcal R_{\varepsilon,B}(t-r)G\,dW_r
 }^2dt\\
 &\qquad=
 \sigma^2\int_0^T(T-s)
 \norm{\mathcal R_{\varepsilon,B}(s)G}_{\HS}^2ds
 \le C_T\varepsilon^4\norm{B}^4
\end{align*}
by Proposition~\ref{prop:noise-transfer}. This proves \eqref{eq:mild-remainder}.

To obtain \eqref{eq:response-mild}, insert \eqref{eq:mild} into the right-hand side. The deterministic term is $\mathcal D_BS(t)u_0$. For the stochastic term, stochastic Fubini is justified by Assumption~\ref{ass:noise}, the semigroup bounds, and boundedness of $B$. It yields
\begin{align*}
 &\sigma\int_0^tS(t-s)B
 \left(\int_0^sS(s-r)G\,dW_r\right)ds\\
 &\qquad=
 \sigma\int_0^t
 \left[
 \int_r^tS(t-s)BS(s-r)ds
 \right]G\,dW_r
 =
 \sigma\int_0^t\mathcal D_BS(t-r)G\,dW_r.
\end{align*}
\end{proof}

\begin{corollary}[Derivative of the field functional]\label{cor:psi-derivative}
For every $w\in\cH$,
\begin{equation}\label{eq:psi-derivative}
 \frac{
 \Psi_w(U^{\varepsilon,B})-\Psi_w(U)
 }{\varepsilon}
 \longrightarrow
 \Psi_w(\dot U[B])
\end{equation}
in $L^2(\Omega)$, uniformly for $B$ in bounded subsets of $\cL(H)$ after division by $\norm{B}$.
\end{corollary}

\begin{proof}
Apply \eqref{eq:functional-cont} to \eqref{eq:mild-remainder}.
\end{proof}

\section{Mean and covariance response}

For $w\in\cH$, define
\begin{equation}\label{eq:mean}
 m_A(w)
 :=
 \int_0^T\ip{w_t}{S(t)u_0}_Hdt.
\end{equation}
Define the adjoint transfer
\begin{equation}\label{eq:g}
 g_r^{A,w}
 :=
 \int_r^TS^*(t-r)w_tdt,
 \qquad 0\le r\le T,
\end{equation}
and the variance
\begin{equation}\label{eq:variance}
 v_A(w)
 :=
 \sigma^2\int_0^T
 \norm{G^*g_r^{A,w}}_U^2dr.
\end{equation}
Stochastic Fubini gives
\begin{equation}\label{eq:Gaussian-law}
 \Psi_w(U)
 \sim
 \N(m_A(w),v_A(w)).
\end{equation}

\begin{theorem}[Mean derivative and remainder]\label{thm:mean-derivative}
For $B\in\cL(H)$,
\begin{equation}\label{eq:Dm}
 D_Am_A(w)[B]
 =
 \int_0^T
 \ip{
 w_t
 }{
 \int_0^tS(t-s)BS(s)u_0ds
 }_Hdt.
\end{equation}
For every $R_w<\infty$,
\begin{align}
 \sup_{\norm{w}_{\cH}\le R_w}
 \abs{D_Am_A(w)[B]}
 &\le C_TR_w\norm{u_0}\norm{B},
 \label{eq:Dm-bound}\\
 \sup_{\norm{w}_{\cH}\le R_w}
 \abs{
 m_{A+\varepsilon B}(w)-m_A(w)
 -\varepsilon D_Am_A(w)[B]
 }
 &\le C_T\varepsilon^2R_w\norm{u_0}\norm{B}^2.
 \label{eq:Dm-remainder}
\end{align}
\end{theorem}

\begin{proof}
Insert \eqref{eq:semigroup-expansion} into the mean and use Cauchy--Schwarz in time. The derivative is \eqref{eq:Dm}; bounds \eqref{eq:Dm-bound}--\eqref{eq:Dm-remainder} follow from Theorem~\ref{thm:semigroup-derivative}.
\end{proof}

\begin{lemma}[Adjoint transfer derivative]\label{lem:g-derivative}
The map $A\mapsto g^{A,w}$ is differentiable in $\cH$, with
\begin{equation}\label{eq:Dg}
 D_Ag_r^{A,w}[B]
 =
 \int_r^T
 \left[
 \int_0^{t-r}
 S^*(s)B^*S^*(t-r-s)ds
 \right]w_tdt.
\end{equation}
Uniformly for bounded $w$-classes,
\begin{align}
 \norm{D_Ag^{A,w}[B]}_{\cH}
 &\le C_T\norm{w}_{\cH}\norm{B},
 \label{eq:Dg-bound}\\
 \norm{
 g^{A+\varepsilon B,w}-g^{A,w}
 -\varepsilon D_Ag^{A,w}[B]
 }_{\cH}
 &\le C_T\varepsilon^2\norm{w}_{\cH}\norm{B}^2.
 \label{eq:Dg-rem}
\end{align}
\end{lemma}

\begin{proof}
Take adjoints in \eqref{eq:DS}:
\[
 \mathcal D_BS(t)^*
 =
 \int_0^tS^*(s)B^*S^*(t-s)ds.
\]
Insert this into \eqref{eq:g}. The stable semigroup bound and Young's inequality for the backward Volterra operator give \eqref{eq:Dg-bound}. Apply the same argument to the adjoint of the remainder in \eqref{eq:semigroup-remainder} to obtain \eqref{eq:Dg-rem}.
\end{proof}

\begin{theorem}[Variance derivative and remainder]\label{thm:variance-derivative}
For $B\in\cL(H)$,
\begin{equation}\label{eq:Dv}
 D_Av_A(w)[B]
 =
 2\sigma^2\int_0^T
 \ip{
 G^*g_r^{A,w}
 }{
 G^*D_Ag_r^{A,w}[B]
 }_Udr.
\end{equation}
For every $R_w<\infty$,
\begin{align}
 \sup_{\norm{w}_{\cH}\le R_w}
 \abs{D_Av_A(w)[B]}
 &\le C_TR_w^2\norm{B},
 \label{eq:Dv-bound}\\
 \sup_{\norm{w}_{\cH}\le R_w}
 \abs{
 v_{A+\varepsilon B}(w)-v_A(w)
 -\varepsilon D_Av_A(w)[B]
 }
 &\le C_T\varepsilon^2R_w^2\norm{B}^2.
 \label{eq:Dv-rem}
\end{align}
\end{theorem}

\begin{proof}
The map $g\mapsto\sigma^2\norm{G^*g}_{L^2(0,T;U)}^2$ is a continuous quadratic form on $\cH$. Apply the product rule together with Lemma~\ref{lem:g-derivative}. The quadratic remainder is the sum of a term involving \eqref{eq:Dg-rem} and the square of the first-order increment. Both are bounded by the right-hand side of \eqref{eq:Dv-rem}.
\end{proof}

\section{Derivative of Gaussian counterfactual laws}

For $v_A(w)>0$, set
\begin{equation}\label{eq:a-standard}
 a_{A,w,z}
 :=
 \frac{z-m_A(w)}{\sqrt{v_A(w)}},
 \qquad
 F_{A,w}(z):=\Ph(a_{A,w,z}).
\end{equation}
Let $\cW\subset\cH$ be bounded, let $\cZ\subset\R$ be compact, and assume
\begin{equation}\label{eq:variance-lower}
 \inf_{w\in\cW}v_A(w)\ge v_->0.
\end{equation}

\begin{theorem}[Uniform CDF derivative]\label{thm:CDF-derivative}
Under \eqref{eq:variance-lower},
\begin{equation}\label{eq:DF}
 D_AF_{A,w}(z)[B]
 =
 \ph(a_{A,w,z})
 \left[
 -\frac{D_Am_A(w)[B]}{\sqrt{v_A(w)}}
 -\frac{z-m_A(w)}{2v_A(w)^{3/2}}
 D_Av_A(w)[B]
 \right].
\end{equation}
Moreover, for sufficiently small $\abs{\varepsilon}$,
\begin{equation}\label{eq:CDF-remainder}
 \sup_{(w,z)\in\cW\times\cZ}
 \abs{
 F_{A+\varepsilon B,w}(z)-F_{A,w}(z)
 -\varepsilon D_AF_{A,w}(z)[B]
 }
 \le C\varepsilon^2\norm{B}^2.
\end{equation}
\end{theorem}

\begin{proof}
Theorems~\ref{thm:mean-derivative} and \ref{thm:variance-derivative} give uniform first-order expansions with quadratic remainders. The lower variance bound and continuity imply that the perturbed variance remains above $v_-/2$ for small $\varepsilon$. The map
\[
 (m,v)\mapsto\Ph((z-m)/\sqrt v)
\]
has uniformly bounded first and second derivatives on the resulting compact set. The Banach-space chain rule gives \eqref{eq:DF}; Taylor's theorem gives \eqref{eq:CDF-remainder}.
\end{proof}

\begin{corollary}[Tail and quantile derivatives]\label{cor:tail-quantile}
The tail probability $R_{A,w}(z)=1-F_{A,w}(z)$ satisfies
\[
 D_AR_{A,w}(z)[B]=-D_AF_{A,w}(z)[B].
\]
For $p\in[p_-,p_+]\Subset(0,1)$, the Gaussian quantile
\begin{equation}\label{eq:quantile}
 Q_{A,w}(p)
 =m_A(w)+\sqrt{v_A(w)}\Ph^{-1}(p)
\end{equation}
has derivative
\begin{equation}\label{eq:DQ}
 D_AQ_{A,w}(p)[B]
 =D_Am_A(w)[B]
 +\frac{\Ph^{-1}(p)}{2\sqrt{v_A(w)}}D_Av_A(w)[B].
\end{equation}
Both derivatives and their quadratic remainders are uniform over the declared compact target classes.
\end{corollary}

\begin{remark}[Where non-degeneracy is needed]
The lower variance bound carries the burden in Theorem~\ref{thm:CDF-derivative} and Corollary~\ref{cor:tail-quantile}. It is not needed for the Wasserstein stability result below, which works directly under the common-noise coupling.
\end{remark}

\section{Wasserstein stability and tangent laws}

\begin{theorem}[Local $W_2$ stability]\label{thm:W2}
Under Assumptions~\ref{ass:semigroup}--\ref{ass:noise},
\begin{equation}\label{eq:W2-bound}
 W_2\left(
 \Law(\Psi_w(U^{\varepsilon,B})),
 \Law(\Psi_w(U))
 \right)
 \le
 C_T\abs{\varepsilon}\norm{B}\norm{w}_{\cH}
\end{equation}
uniformly over $B\in\mathbb B_R$. In addition,
\begin{equation}\label{eq:tangent-law}
 W_2\left(
 \Law\left(
 \frac{\Psi_w(U^{\varepsilon,B})-\Psi_w(U)}{\varepsilon}
 \right),
 \Law(\Psi_w(\dot U[B]))
 \right)
 \longrightarrow0.
\end{equation}
\end{theorem}

\begin{proof}
Use the coupling in which the perturbed and unperturbed fields share the same Wiener process. By the definition of $W_2$, \eqref{eq:functional-cont}, and the first-order semigroup bound,
\begin{align*}
 W_2^2
 &\le
 \E\abs{\Psi_w(U^{\varepsilon,B}-U)}^2\\
 &\le
 \norm{w}_{\cH}^2
 \norm{U^{\varepsilon,B}-U}_{\cX}^2
 \le C_T\varepsilon^2\norm{B}^2\norm{w}_{\cH}^2.
\end{align*}
This proves \eqref{eq:W2-bound}. For \eqref{eq:tangent-law}, use the same coupling and Corollary~\ref{cor:psi-derivative}.
\end{proof}

\begin{remark}[Gaussian formula]
For the one-dimensional target laws,
\begin{equation}\label{eq:Gaussian-W2}
 W_2^2(\N(m_1,v_1),\N(m_0,v_0))
 =(m_1-m_0)^2+(\sqrt{v_1}-\sqrt{v_0})^2.
\end{equation}
The coupling proof is preferable for stability because it does not require a lower variance bound and preserves the structural shared-noise interpretation from WP1.
\end{remark}

\section{The unbounded scalar-rescaling intervention}

Let $A_0$ generate an exponentially stable analytic semigroup. For
\begin{equation}\label{eq:a-range}
 a\in[a_-,a_+]\Subset(0,\infty),
\end{equation}
write
\begin{equation}\label{eq:Sa}
 S_a(t):=e^{taA_0}.
\end{equation}

\begin{assumption}[Uniform analytic bounds]\label{ass:analytic}
For $j=1,2$,
\begin{equation}\label{eq:analytic-bound}
 \sup_{a\in[a_-,a_+]}
 \sup_{t>0}
 \norm{t^jA_0^jS_a(t)}<\infty.
\end{equation}
This holds for the negative self-adjoint generator in the paper's diagonal core and, more generally, for bounded analytic semigroups uniformly over compact positive rescalings.
\end{assumption}

\begin{theorem}[Sensitivity of scalar generator rescaling]\label{thm:scalar-scaling}
Under Assumption~\ref{ass:analytic},
\begin{equation}\label{eq:scalar-derivatives}
 \partial_aS_a(t)=tA_0S_a(t),
 \qquad
 \partial_a^2S_a(t)=t^2A_0^2S_a(t).
\end{equation}
Uniformly for $a\in[a_-,a_+]$ and $t\in[0,T]$,
\begin{equation}\label{eq:scalar-rem}
 \norm{S_{a+h}(t)-S_a(t)-h\,tA_0S_a(t)}
 \le Ch^2.
\end{equation}
If the corresponding noise-transfer derivatives are square-integrable in Hilbert-Schmidt norm, all conclusions of Theorems~\ref{thm:mild-derivative}--\ref{thm:W2} hold with the scalar direction
\begin{equation}\label{eq:dotSa}
 \dot S_a(t)=tA_0S_a(t).
\end{equation}
\end{theorem}

\begin{proof}
For $t>0$, analytic semigroup calculus gives \eqref{eq:scalar-derivatives}. Taylor's formula in the scalar parameter yields
\[
 S_{a+h}(t)-S_a(t)-h\,tA_0S_a(t)
 =h^2\int_0^1(1-r)t^2A_0^2S_{a+rh}(t)dr.
\]
The uniform bound \eqref{eq:analytic-bound} proves \eqref{eq:scalar-rem}; at $t=0$ both sides vanish. The remaining claims follow by repeating the deterministic and stochastic arguments above with \eqref{eq:dotSa} and the stated noise-transfer integrability.
\end{proof}

\begin{corollary}[Diagonal formulas used by WP4]\label{cor:diagonal}
Suppose $A_0e_k=-\lambda_ke_k$ and $G e_k=q_ke_k$. Then
\begin{align}
 \dot m_a(w)
 &=\int_0^T\ip{w_t}{tA_0S_a(t)u_0}dt,
 \label{eq:dotm-diag}\\
 (\dot K_aw)_k(r)
 &=-q_k\int_r^T\lambda_k(t-r)e^{-a\lambda_k(t-r)}w_k(t)dt,
 \label{eq:dotK-diag}\\
 \dot V_a(w)
 &=2\ip{K_aw}{\dot K_aw}_{L^2(0,T;H)}.
 \label{eq:dotV-diag}
\end{align}
The resulting CDF derivative is exactly the intervention derivative used in the WP4 process theorem.
\end{corollary}

\begin{remark}[Spatially varying diffusivity]
A direction such as $B=\nabla\!\cdot(b\nabla)$ is generally unbounded and may change the generator domain. It is not covered by Theorem~\ref{thm:semigroup-derivative}. A rigorous extension requires quadratic-form perturbations, relative boundedness, or an extrapolation-space framework. WP6 does not claim that extension as proved.
\end{remark}

\section{Causal response geometry}

For fixed $(w,z)$, Theorem~\ref{thm:CDF-derivative} defines the bounded linear functional
\begin{equation}\label{eq:covector}
 \mathfrak r_{A,w,z}:\cL(H)\to\R,
 \qquad
 \mathfrak r_{A,w,z}(B):=D_AF_{A,w}(z)[B].
\end{equation}

\begin{remark}[Covector versus gradient]
In the operator norm, \eqref{eq:covector} is a response covector. The Banach space $\cL(H)$ has no canonical Hilbert inner product, so there is no intrinsic gradient. A gradient appears only after an admissible perturbation geometry is chosen.
\end{remark}

Restrict to the Hilbert space $\HS(H)$ with
\begin{equation}\label{eq:HS-inner}
 \ip{C}{B}_{\HS}:=\tr(C^*B).
\end{equation}
For $a,b\in H$, write
\begin{equation}\label{eq:rank-one}
 (a\otimes b)x:=\ip{b}{x}a.
\end{equation}
Then
\begin{equation}\label{eq:rank-identity}
 \ip{a\otimes b}{B}_{\HS}=\ip{a}{Bb}_H.
\end{equation}
Let $\CG:=GG^*$.

\begin{theorem}[Hilbert-Schmidt mean and variance gradients]\label{thm:HS-gradients}
For $B\in\HS(H)$,
\begin{align}
 D_Am_A(w)[B]
 &=\ip{\nabla_Am_A(w)}{B}_{\HS},
 \label{eq:mean-gradient-id}\\
 D_Av_A(w)[B]
 &=\ip{\nabla_Av_A(w)}{B}_{\HS},
 \label{eq:variance-gradient-id}
\end{align}
where
\begin{equation}\label{eq:mean-gradient}
 \nabla_Am_A(w)
 =
 \int_0^T\int_0^t
 [S^*(t-s)w_t]\otimes[S(s)u_0]
 \,ds\,dt
\end{equation}
and
\begin{equation}\label{eq:variance-gradient}
 \nabla_Av_A(w)
 =
 2\sigma^2
 \int_0^T\int_r^T\int_0^{t-r}
 [S^*(t-r-s)w_t]
 \otimes
 [S(s)\CG g_r^{A,w}]
 \,ds\,dt\,dr.
\end{equation}
The integrals are Bochner integrals in $\HS(H)$.
\end{theorem}

\begin{proof}
For the mean, write the integrand of \eqref{eq:Dm} as
\[
 \ip{S^*(t-s)w_t}{BS(s)u_0}
\]
and use \eqref{eq:rank-identity}. Stable semigroup bounds make \eqref{eq:mean-gradient} Bochner integrable.

For the variance, insert \eqref{eq:Dg} into \eqref{eq:Dv}. With
$a=S^*(t-r-s)w_t$ and $b=S(s)\CG g_r^{A,w}$,
\begin{align*}
 \ip{\CG g_r^{A,w}}{S^*(s)B^*a}
 &=\ip{S(s)\CG g_r^{A,w}}{B^*a}
 =\ip{a}{Bb}.
\end{align*}
Use \eqref{eq:rank-identity} and Fubini. Integrability follows from exponential stability, boundedness of $\CG$, and Cauchy--Schwarz in the $r$ and $t$ variables.
\end{proof}

\begin{corollary}[CDF gradient and optimal local direction]\label{cor:CDF-gradient}
Under the assumptions of Theorem~\ref{thm:CDF-derivative},
\begin{equation}\label{eq:CDF-gradient}
 \nabla_AF_{A,w}(z)
 =
 \ph(a_{A,w,z})
 \left[
 -\frac{\nabla_Am_A(w)}{\sqrt{v_A(w)}}
 -\frac{z-m_A(w)}{2v_A(w)^{3/2}}
 \nabla_Av_A(w)
 \right].
\end{equation}
If this gradient is non-zero, then
\begin{equation}\label{eq:optimal-direction}
 \sup_{\norm{B}_{\HS}\le1}
 D_AF_{A,w}(z)[B]
 =
 \norm{\nabla_AF_{A,w}(z)}_{\HS},
\end{equation}
attained at
\begin{equation}\label{eq:Bstar}
 B^*
 =
 \frac{\nabla_AF_{A,w}(z)}
 {\norm{\nabla_AF_{A,w}(z)}_{\HS}}.
\end{equation}
For a closed Hilbert subspace of admissible directions, replace the gradient by its orthogonal projection onto that subspace.
\end{corollary}

\begin{proof}
Equation \eqref{eq:CDF-gradient} follows from \eqref{eq:DF} and Theorem~\ref{thm:HS-gradients}. Equations \eqref{eq:optimal-direction}--\eqref{eq:Bstar} are the Riesz representation theorem and Cauchy--Schwarz inequality in $\HS(H)$.
\end{proof}

\begin{remark}[Structural constraints]
If perturbations must be self-adjoint, use the symmetric projection of \eqref{eq:CDF-gradient}. If they must preserve mass, locality, sparsity, or a finite parametric operator family, project onto the corresponding admissible tangent space. The response geometry is meaningful only relative to such a declared intervention class.
\end{remark}

\section{Finite-dimensional validation}

The numerical experiment uses a six-dimensional stable diagonal generator and a symmetric off-diagonal perturbation that does not commute with it. The state and adjoint transfer are solved by high-accuracy ODE integration. Their sensitivity equations are
\begin{align}
 x'(t)&=Ax(t),& y'(t)&=Ay(t)+Bx(t),
 \label{eq:forward-sens}\\
 g'(r)&=-w(r)-A^*g(r),& h'(r)&=-B^*g(r)-A^*h(r),
 \label{eq:backward-sens}
\end{align}
with $x(0)=u_0$, $y(0)=0$, $g(T)=h(T)=0$. Then
\begin{align}
 \dot m&=\int_0^T\ip{w(t)}{y(t)}dt,\
 \dot v&=2\sigma^2\int_0^T\ip{G^*g(r)}{G^*h(r)}dr.
\end{align}

Table~\ref{tab:numerical} shows the central finite-difference CDF derivative error and the first-order law remainder.

\begin{table}[ht]
\centering
\small
\begin{tabular}{@{}rrrr@{}}
\toprule
$\abs{\varepsilon}$ & Central derivative error & First-order CDF remainder & Exact $W_2$\\
\midrule
0.100 & $4.27\times10^{-6}$ & $1.16\times10^{-4}$ & $1.28\times10^{-3}$\\
0.050 & $1.07\times10^{-6}$ & $2.92\times10^{-5}$ & $6.25\times10^{-4}$\\
0.020 & $1.71\times10^{-7}$ & $4.67\times10^{-6}$ & $2.46\times10^{-4}$\\
0.010 & $4.27\times10^{-8}$ & $1.17\times10^{-6}$ & $1.22\times10^{-4}$\\
0.005 & $1.07\times10^{-8}$ & $2.92\times10^{-7}$ & $6.08\times10^{-5}$\\
0.002 & $1.71\times10^{-9}$ & $4.67\times10^{-8}$ & $2.43\times10^{-5}$\\
\bottomrule
\end{tabular}
\caption{Non-commuting operator-sensitivity validation.}
\label{tab:numerical}
\end{table}

\begin{figure}[H]
\centering
\includegraphics[width=.76\textwidth]{wp6_finite_difference_derivative.png}
\caption{Forward and central finite-difference errors against perturbation size. The central difference follows the expected second-order slope.}
\label{fig:fd}
\end{figure}

\begin{figure}[H]
\centering
\includegraphics[width=.76\textwidth]{wp6_cdf_linearisation_remainder.png}
\caption{First-order CDF linearisation remainder. The empirical error is quadratic in $\abs{\varepsilon}$.}
\label{fig:cdf-rem}
\end{figure}

\begin{figure}[H]
\centering
\includegraphics[width=.76\textwidth]{wp6_wasserstein_stability.png}
\caption{Exact Gaussian $W_2$ distance and the common-noise Lipschitz envelope.}
\label{fig:w2}
\end{figure}

\section{Completion statement and WP7 interface}

The package declares the operator topology, proves all derivative and interchange steps, gives explicit law-level derivatives, and verifies them numerically. The interface with WP7 is:
\begin{enumerate}[label=(\roman*)]
\item bounded operator-approximation errors can be propagated through Theorems~\ref{thm:semigroup-derivative}, \ref{thm:CDF-derivative}, and \ref{thm:W2};
\item scalar diffusivity approximation uses Theorem~\ref{thm:scalar-scaling};
\item a certified second-order remainder allows WP7 to separate linear operator bias from higher-order error;
\item genuinely unbounded, domain-changing approximation directions require a separate form-perturbation extension and are not hidden inside the present theorem.
\end{enumerate}

\begin{thebibliography}{9}

\bibitem{DaPratoZabczyk2014}
G. Da Prato and J. Zabczyk.
\newblock \emph{Stochastic Equations in Infinite Dimensions}, second edition.
\newblock Cambridge University Press, 2014.

\bibitem{EngelNagel2000}
K.-J. Engel and R. Nagel.
\newblock \emph{One-Parameter Semigroups for Linear Evolution Equations}.
\newblock Springer, 2000.

\bibitem{Marinelli2020}
C. Marinelli.
\newblock On the differentiability of solutions to singularly perturbed SPDEs.
\newblock arXiv:2012.15338, 2020.

\bibitem{Pazy1983}
A. Pazy.
\newblock \emph{Semigroups of Linear Operators and Applications to Partial Differential Equations}.
\newblock Springer, 1983.

\bibitem{Villani2009}
C. Villani.
\newblock \emph{Optimal Transport: Old and New}.
\newblock Springer, 2009.

\end{thebibliography}

\end{document}
